\section{Demonstration and Findings}

The demonstration evaluates the resource interface rather than ranking
tokenizer systems. It combines two public SID export paths with controlled rows
that calibrate the diagnostics. GRID supplies the open RQ-KMeans-style path;
ReSID supplies the GAOQ line and processed Musical artifacts. Sanity tokenizers
and stressors are included only when they clarify what a diagnostic can and
cannot identify.

\paragraph{Case-study protocol.}
The paper-facing case study uses the same 23,742 Musical items for both rows.
The GRID row is a controlled feature-text export: ReSID processed feature text
is embedded and passed through the official GRID-style residual k-means module.
The ReSID row is a bounded FAMAE-to-GAOQ export on the same item universe. This
modest protocol gives reviewers a same-item diagnostic contrast while avoiding
the stronger, unsupported claim that the paper reproduces every training choice
of TIGER, GRID, or the full ReSID pipeline.

\begin{table}[t]
\caption{Same-item Musical artifact profiles on 23,742 items. Higher unique
full codes, D3 L1 co-occurrence prefix recall, and D4 tail unique-SID ratio are
better; lower D2 full-collision rate is better. D5a prefixes are level-wise
prefix counts, not latency. A dagger marks structural one-full-code-per-item
rows where D2=0 and D4 tail=1.0 are forced by the export design.}
\label{tab:musical-diagnostic}
\footnotesize
\setlength{\tabcolsep}{1.0pt}
\begin{tabular}{@{}>{\raggedright\arraybackslash}p{0.24\columnwidth}rrrr>{\raggedright\arraybackslash}p{0.19\columnwidth}@{}}
\toprule
System & Unique & D2 full & D3 L1 & D4 tail & D5a prefixes \\
\midrule
\textbf{GRID ft} (3s) & 3,857$\pm$112 & .9759$\pm$.0009 & .0519$\pm$.0037 & .377$\pm$.008 & 64/3.4--3.7k/3.7--4.0k \\
\textbf{ReSID}\textsuperscript{\dag} & 23,742 & .0000 & .1535 & 1.0000 & 32/1280/23.7k \\
\emph{Cat-prefix}\textsuperscript{\dag} & 23,742 & .0000 & .4470 & 1.0000 & 30/83/313/23.7k \\
\emph{Hash-collide} & 256 & 1.0000 & .0037 & .0322 & 256/256/256/256 \\
\emph{Pop-balanced} & 22,707 & .0860 & .3026 & .9619 & 4/1024/22.7k \\
\bottomrule
\end{tabular}
\end{table}

Table~\ref{tab:musical-diagnostic} compares two exports on the same item set.
The GRID row is intentionally labeled as feature-text: it uses processed
feature text from the ReSID artifact path and the official residual
quantization module, not a faithful raw-text TIGER reproduction. Under that
controlled setup, the row exposes high collision pressure and limited tail
capacity. The ReSID/GAOQ row exports one full code per item in this bounded
run, yielding zero full-code collisions and higher D3-L1 collaborative prefix
recovery. The daggered rows make the structural floor explicit: collision-free
full codes are not by themselves evidence of better recommendation quality.
Across seeds 42--44, the GRID duplicate-\sid rate remains 0.8327--0.8421 and
full-collision rate remains 0.9751--0.9769, so the collision pressure is not a
single-seed artifact. The comparison is also architecturally biased: GRID is
capacity-constrained, while the bounded ReSID export reaches an item-unique
leaf level. A same-capacity ablation is future work, not a hidden claim.

\begin{table}[t]
\caption{Controlled stressors isolate diagnostic mechanisms. Each row activates
one failure mode and asks whether the corresponding diagnostic separates the
two conditions. Full numeric tables are shipped in the artifact.}
\label{tab:controller-stressors}
\scriptsize
\setlength{\tabcolsep}{2pt}
\begin{tabular}{@{}>{\raggedright\arraybackslash}p{0.20\columnwidth}>{\raggedright\arraybackslash}p{0.13\columnwidth}>{\raggedright\arraybackslash}p{0.24\columnwidth}>{\raggedright\arraybackslash}p{0.17\columnwidth}>{\raggedright\arraybackslash}p{0.17\columnwidth}@{}}
\toprule
Stressor & D & Test & Baseline & Under stress \\
\midrule
Qualified collision & D2b & co-occur $\ne$ hash collisions & hash: 1.19x risk & co-occur: 3.86x risk \\
Capacity budget & D1/D4 & tight codebooks skew capacity & head: 1.000 unique & tail: 0.028 unique \\
Variable depth & D5a & active cost $\ne$ max-depth view & max-depth: 12,010 & active: 7,914 \\
\bottomrule
\end{tabular}
\end{table}

\paragraph{Diagnostic findings.}
Tables~\ref{tab:musical-diagnostic} and~\ref{tab:controller-stressors} support
three artifact-level findings. First, collision pressure is stable in the
controlled GRID row, while the same interface ingests a collision-free ReSID
mapping. The qualified-collision stressor then shows why raw collision volume
is insufficient: collisions that overlap user co-occurrence have a 3.86x lift,
while synthetic hash collisions have 1.19x. Second, collision-free addresses
and collaborative-prefix alignment are different objectives. ReSID has no full
collisions and higher D3-L1 than the GRID row, while a category-prefix control
in the repository recovers more level-1 co-occurrence neighbors without being a
learned tokenizer. Third, D4 and D5a separate capacity allocation from prefix
structure: in the capacity stressor, head uniqueness reaches 1.000 while tail
uniqueness remains 0.028, and variable-depth stressors show that active prefix
cost (7,914 realized prefixes) need not match a fixed maximum-depth view
(12,010 prefixes).

Repository tables extend these checks without changing the paper claims:
sanity rows calibrate D1--D5a, GRID scale rows test larger item counts, DACT
rows exercise optional D6 churn, and MovieLens rows check non-Amazon schema
portability. Table~\ref{tab:musical-diagnostic} remains the only method-facing
case study in the main paper.
